% =============================================================================
\chapter[Your draft, read against the measurements]{Your draft, read against\\ the measurements}
\label{ch:review}

I read your draft of 16 September as a referee would and against everything in the previous chapters, and
its theory was re-implemented from the definitions and searched for counterexamples. None was found: the
mathematics looks right, and every finding below is about what the text says. The numbers printed in your
draft were not audited in that pass. Each item gives where it sits, why it matters, and the fix.

\begin{ideia}[The review in two sentences]
Keep the object, the fixpoint definition and the reporting discipline as they are. Fix the positioning
against the paper you build on, the conditions the text drops, and prose that claims more than its table.
\end{ideia}

\begin{antes}
\antesq[espaco]{Which of your structural results does the foundation paper state to be false for
MPDAGs in general?}
\antesq[distancia-em-saltos]{Is there a formula for $d(G, H)$ on your space that needs no search?}
\end{antes}

\section{What the draft gets right}

\textbf{The object, in the right geometry.} Definition~1 by the fixpoint test is a real contribution:
Appendix~A.1 shows with minimal witnesses (the collider on a three-node chain, $K_4$ minus an edge with
one claim) why the two obvious filters are unsound, and covers are computed from model inclusion rather
than from a conjectured characterisation.

\textbf{The clinician sentence} in \S4.4 is the best-written sentence of the draft, and it makes the
object contestable by a domain expert. It belongs on page one.

\textbf{The reporting discipline} of \S5.4 is already the one our protocol uses: statuses on a frozen
instance list rather than filters, censoring rather than dropping, no numeric radius for an unreachable
failure, the two corruption arms never merged. \S6 states the observables-only invariant of
Chapter~\ref{ch:vocabulary} in one sentence. And Appendix~B's E3 is designed to refute rather than to
confirm («agreement proves nothing; only disagreement would be a theorem»), which is the right
epistemics for a verified-not-proved property.

\begin{medido}[Your theory, re-implemented from Definitions 1 to 3]
The space, the order, the covering relation and $d$ were rebuilt using only \texttt{MPDAG},
\texttt{apply\_orientations} and \texttt{enumerate\_dag\_extensions} from \texttt{src/\allowbreak bkrobust/\allowbreak demo/\allowbreak },
never the search or SAT code. Scale: 1,375 spaces, about 61,000 knowledge states, 9,648,249 ordered
pairs, 46,698 covering pairs, 188,116 anti-exchange triples (39,402 of them on states whose undirected
component is not chordal). Retraction-only radius equal to the brute-force radius on 8,161 of 8,161
instances; 599 of 599 nearest failing states reachable by retractions; the SAT clause set has 1,858
models for 1,858 space elements. \textbf{No counterexample to any checkable statement.} The check shares
your Meek and extension primitives, so a bug there would be invisible to it. From my review notes of 16
September; not in the branch.
\end{medido}

% -----------------------------------------------------------------------------
\begin{antes}
\antesq[espaco]{Where in the draft is the anti-exchange property stated?}
\antesq[gac-backdoor]{Which paragraph of the draft computes radii for sets other than the optimal one?}
\end{antes}

\section{The theory: what the text says}

\figref{fig:review-chain} is the dependency chain as the appendix proves it. Three results are
unconditional; four inherit the open case of the anti-exchange property; and the property itself is
the one statement the reader is never shown.

\begin{figure}[htbp]
\centering
\begin{tikzpicture}[font=\scriptsize\sffamily]
\node[cmapr] (ae) at (4.4,0) {anti-exchange:\\ stated nowhere, 4 × ??};
\node[cmap] (l2) at (4.4,-1.0) {Lemma 2: a cover\\ adds one orientation};
\node[cmap] (ps) at (4.4,-2.0) {Property 1: upper\\ semimodularity};
\node[cmap] (l3) at (4.4,-3.0) {Lemma 3: the join\\ is no farther};
\node[cmap] (p1) at (4.4,-4.0) {Prop.~1: monotone\\ reachability};
\node[cmapg] (l1) at (0,-1.0) {Lemma 1: order\\ identification};
\node[cmapg] (t2) at (0,-2.0) {Theorem 2:\\ join-semilattice};
\node[cmapg] (t1) at (0,-4.0) {Theorem 1: failure\\ upward-closed};
\node[cmapw] (tgh) at (9.0,0) {Taeb--Guo--Henckel, Prop.~9:\\ all MPDAGs on $V$, not graded};
\node[cmap] (gr) at (9.0,-1.0) {graded, rank $|\Kset_G|$};
\node[cmap] (cf) at (9.0,-2.0) {closed form\\ $d = |\Kset_G \triangle \Kset_H|$};
\node[cmap] (se) at (9.0,-4.0) {retraction-only search;\\ E1 counts equal hops};
\node[cmapc] (drop) at (9.0,-3.0) {condition dropped in abstract,\\ contributions, conclusion, \S5.3};
\draw[cedge] (ae) -- node[right]{via Edelman--Jamison} (l2);
\draw[cedge, densely dashed] (l1) -- node[above]{not cited there} (l2);
\draw[cedge] (l2) -- (ps);
\draw[cedge] (t2) -- (ps);
\draw[cedge] (ps) -- (l3);
\draw[cedge] (l3) -- (p1);
\draw[cedge] (t1) -- (p1);
\draw[cedge] (l2) -- (gr);
\draw[cedge] (l2) -- (cf);
\draw[cedge] (p1) -- (se);
\draw[cedge] (se) -- node[right]{stated as unconditional} (drop);
\draw[cedge, densely dashed, vermesc] (tgh) -- node[right]{skeleton varies there} (gr);
\end{tikzpicture}
\caption{Green results hold unconditionally; every white box inherits the anti-exchange property,
which the draft never states. The foundation paper's non-gradedness is about a larger poset, and the
draft does not yet say so.}
\label{fig:review-chain}
\end{figure}

\subsection*{1. The foundation paper states non-gradedness for a larger poset}

\emph{Where.} Appendix~A calls $\rho(G) = |\mathrm{dir}(G)| - |\mathrm{dir}(\Chat)|$ a rank function «so
the order is graded», Property~1 is upper semimodularity, and Theorem~2 makes the space a
join-semilattice. \emph{Why.} Proposition~9 of Taeb, Guo and Henckel states that the model-oriented poset
of MPDAGs on $|V| \ge 4$ vertices is not graded, and the paragraph after it says the poset is neither upper
nor lower semimodular, has no semilattice structure, and has no explicit graphical characterisation of
its covering relation; a referee who knows that paper reads Appendix~A as a contradiction. \emph{Fix.}
Say why both hold, in one paragraph of \S5.1 and one sentence of Appendix~A.

\begin{fonte}[Taeb, Guo and Henckel, arXiv:2511.10625v3, \S4.5 and Appendix D.3]
Their poset contains every MPDAG on $V$, from the empty graph to the complete undirected one, and covers
can change adjacencies: in their Example~2 two neighbours of the same graph differ from it by one edge
and by $n-1$ edges, and the proof of Proposition~9 compares two maximal chains of lengths 3 and 5 in
which a covering step removes an edge. On a common skeleton their Proposition~8 reduces the order to $[G] \subseteq [H]$,
which is yours, and their pseudo-rank (Definition~7, directed edges plus twice the undirected ones) equals
$2m - |\mathrm{dir}(G)|$ on a skeleton with $m$ edges: your rank, negated and shifted.
\end{fonte}

The paragraph needs five sentences, and it turns the largest liability of the draft into its clearest
positioning statement.

\begin{sonorio}[A paragraph for \S5.1]
These results do not contradict Proposition 9 of Taeb, Guo and Henckel (2025), which shows that the
model-oriented poset of all MPDAGs on four or more vertices is not graded, is not a semilattice and is
not semimodular. Their poset lets a covering step change the skeleton. The space studied here fixes the
skeleton and the v-structures of $\widehat{C}$, so only orientations vary. On it the join exists
unconditionally (Theorem 2), and gradedness and upper semimodularity hold under Conjecture 1, with their
pseudo-rank reducing to $2m - |\mathrm{dir}(G)|$.
\end{sonorio}

\subsection*{2. The anti-exchange property is never stated}

\emph{Where.} Four \verb|\Cref{prop:antiexchange}| calls (\S7, the proof of Lemma~2, the heading of
Appendix~A.3, the E3 paragraph of Appendix~B) point at a label that exists nowhere, so the compiled PDF
prints «??» four times, including the heading «A.3 Status of ??». \emph{Why.} The reproducibility
statement promises that the single unproved property is isolated in A.3, and the reader finds a section
about a statement they have never seen. \emph{Fix.} One environment, which also repairs the four
references.

\begin{noprato}[The statement, as \texttt{THEOREMS.md} §4c gives it]
Add to Appendix~A a \texttt{conjecture} environment (your macros define it) carrying
\texttt{\textbackslash label\{prop:antiexchange\}}: \emph{for a closed set $S$ of orientations and
distinct orientations $x, y \notin S$, if $y \in \mathrm{cl}(S \cup \{x\})$ then $x \notin \mathrm{cl}(S
\cup \{y\})$}. Its reading in words: no two distinct orientations are perfectly correlated across the
represented DAGs. Calling it a conjecture matches A.3, where Case~B is open.
\end{noprato}

\subsection*{3. Your theorems make the distance a closed form}

\emph{Where.} Definition~2 defines $d$ as a shortest path, and Appendix~C's «Against a naive count»
argues that counting the assertions that differ understates $d$. \emph{Why.} Lemma~2 and Theorem~2 give
$d(G, H) = |\Kset_G \triangle \Kset_H|$: each covering step changes $\Kset$ by exactly one orientation,
so $d \ge |\Kset_G \triangle \Kset_H|$; and $J = G \vee H$ has $\Kset_J = \Kset_G \cap \Kset_H$, so the
two maximal chains up to $J$ have $|\Kset_G \setminus \Kset_H|$ and $|\Kset_H \setminus \Kset_G|$ steps
and $d \le |\Kset_G \triangle \Kset_H|$. The re-implementation found the identity on all 3,047,289
ordered pairs on which it computed $d$ (a subset of the 9,648,249 above), with no exception and no
unreachable pair, so the naive-count paragraph argues against a count the theory makes equal. That
paragraph reports a Spearman correlation of 0.73 to 0.77 between $d$ and its naive count
(\texttt{90\_appendix.tex}); the other natural reading, a flip counted once, does not give 0.73 to
0.77 either. \emph{Fix.} State the closed form as a
corollary of Lemma~2 (your E2 paragraph already writes $|\Kset_{\Gz} \triangle \Kset_G|$), and either
drop the naive-count paragraph or name exactly which count it compares.

\prompt[conceito=distancia-em-saltos,dificuldade=0.75]{State the closed form of $d$ on the space of
knowledge states, the two results it follows from, and which inequality each gives.}%
{$d(G, H) = |\Kset_G \triangle \Kset_H|$. Lemma 2 (a cover changes $\Kset$ by exactly one orientation)
gives $d \ge |\Kset_G \triangle \Kset_H|$, since each step moves the symmetric difference by at most
one. Theorem 2 gives the join $J$ with $\Kset_J = \Kset_G \cap \Kset_H$, and with Lemma 2 again the
maximal chains from $G$ and from $H$ up to $J$ have $|\Kset_G \setminus \Kset_H|$ and
$|\Kset_H \setminus \Kset_G|$ steps, so $d \le |\Kset_G \triangle \Kset_H|$. It inherits the
anti-exchange condition through Lemma 2.}

\subsection*{4. Proposition 1 says more than its proof, and the condition disappears where it is read}

\emph{Where.} Proposition~1 reads «The nearest failing state is reachable\dots», while its proof
replaces the nearest failing state $G$ by $J = G \vee \Gz$. \S5.1 promises «We state this caveat
wherever a radius is reported», and the abstract («enough order structure for the radius to be computed
without enumerating it»), the contributions («so the search can be restricted to retractions»), the
conclusion («which confines the search to retractions»; «exact radii») and \S5.3 («guaranteed to return
identical radii») state it without the condition. \emph{Why.} The proof delivers some failing state at
distance at most $\rval$ reachable by retractions only, which is what the algorithm needs; the two proved results
named in the contributions do not confine the search, the conditional proposition does, and Appendix~B
itself says all three encodings are upper bounds. \emph{Fix.} Write «some nearest failing state», and
add «under Conjecture~1» in those four places. The strong reading survived on 599 of 599 nearest
failing states, which is evidence, not proof.

\subsection*{5. The Edelman--Jamison appeal lacks its hypotheses}

\emph{Where.} The proof of Lemma~2. \emph{Why.} Their theorem is about a closure operator defined on
every subset of a finite ground set, while here $\mathrm{cl}$ is undefined on inconsistent orientation
sets (937 of 4,137 two-orientation subsets, 22.7\,\%, over 150 spaces), the full ground set was closed in
none of those 150 spaces, and the identification of your order with reverse inclusion of $\Kset_G$ is
Lemma~1, not cited at that point. \emph{Fix.} Prove Lemma~2 directly from the anti-exchange statement, or
say how $\mathrm{cl}$ is extended to a total operator and cite Lemma~1 there.

\subsection*{6. Two adjustment criteria, and where they part}

\emph{Where.} \S2 cites the generalized adjustment criterion; the fast oracle in
\texttt{src/\allowbreak bkrobust/\allowbreak mpdag\_\allowbreak criterion/\allowbreak criterion.py} decides the back-door strengthening, as its docstring
says; Appendix~C's «Robust selection» computes $\rval$ for all fourteen valid sets. \emph{Why.}
\texttt{THEOREMS.md} §14 proves the two criteria coincide on the optimal set of the graph being
judged. Your radius judges one fixed set, $\Ostar(\Gz)$, in other knowledge states, where it is not
their optimal set, so the proof does not cover it; your own BFS check in §14 found identical radii on
5,304 instances, and the ledger audit found the back-door function rejecting 4 of 1,169 committed sets
at the truth (Chapter~\ref{ch:table4}). The main text is therefore very likely unaffected, by
measurement rather than by proof. On arbitrary sets the two disagree on 23,034 of 372,430 cases
(6.2\,\%) in the re-implementation, none of them at the optimal set of the judged graph, and \texttt{THEOREMS.md} §14 records radii that differ on 29.14\,\% of 16,926
arbitrary-set instances, so «Robust selection» is exactly the regime where the choice shows. The
criterion for maximally oriented graphs is introduced in \texttt{perkovic2017cpdag}, already in your
bibliography, while \S1 and \S2 cite \texttt{perkovic2018complete}, whose title restricts it to ancestral
graphs. \emph{Fix.} Name the criterion the code decides, say which one the fourteen sets were scored
with, and move the MPDAG attribution to \texttt{perkovic2017cpdag}.

\prompt[conceito=gac-backdoor,dificuldade=0.7]{Why is the main text of the draft very likely
unaffected by the difference between the GAC and the back-door criterion, why is that not a proof,
and which paragraph is clearly affected?}%
{\texttt{THEOREMS.md} §14 proves the two criteria coincide on the optimal set of the graph being
judged. The main text judges the fixed set $Z = \Ostar(\Gz)$ in other states, where it is an arbitrary
set, so the proof does not reach it; the evidence is measured instead (identical radii on 5,304
instances in §14's BFS check), and the ledger shows the criteria can part there (4 of 1,169 committed
sets rejected by back-door at the truth). Appendix C's «Robust selection» computes radii for every valid adjustment
set, which is the regime where the two criteria part (6.2\,\% of 372,430 cases), so it must say which
criterion it used.}

\prompt[conceito=espaco,dificuldade=0.7]{What does Proposition 9 of Taeb, Guo and Henckel state, and why
does it not contradict the rank function of your Appendix A?}%
{The model-oriented poset of all MPDAGs on $|V| \ge 4$ vertices is not graded (and, by the paragraph
after it, not semimodular and without a characterisation of its covers). Their poset lets covers change
the skeleton; the space of knowledge states fixes the skeleton and the v-structures of $\Chat$, so only
orientations vary. There their pseudo-rank equals $2m - |\mathrm{dir}(G)|$ and the draft's rank is its
negation shifted, so gradedness can hold (under the anti-exchange condition) without contradicting
Proposition 9.}

% -----------------------------------------------------------------------------
\begin{antes}
\antesq[frases-proibidas]{Which cells of your Table~1 sit under the sentence «a strong and consistent
predictor»?}
\antesq[rhostar-raio-de-ruptura]{Where does the word «breakdown» come from in econometrics?}
\end{antes}

\section{The results section and the manuscript}

\subsection*{7. The prose contradicts the table on the same page}

\emph{Where.} \S6.1, «Ranking quality» and «Over-claiming does not pay», printed on page~9 with Table~1.
\emph{Why.} The sentences say more than the table, as the box shows. \emph{Fix.} Uncomment the
«Limits of the present evidence» paragraph of \texttt{07\_discussion.tex}, whose claim is the right one
(«the radius shares significance with the baselines in most strata \dots\ not that it dominates them»),
and restrict «more claims, less robustness» to the tiered arm.

\begin{armadilha}[The sentence against the cells printed under it]
«A strong and consistent predictor»: $\tau_b(\rval) = -0.091\ [-0.204, 0.030]$ in \emph{flip, cov 1.0, bw
0.00} (a cell your caption marks unstable), and the confirmed null $0.032\ [-0.112, 0.168]$ in \emph{flip,
cov 0.5, bw 0.25} at 1,000 draws. «Standard baseline metrics fail»: in
\emph{flip, cov 1.0, bw 0.25}, SHD scores 0.582 and $|\Kset|$ scores 0.557 against the radius's 0.355.
«More claims mean less robustness»: $\tau_b(|\Kset|)$ is positive in five of the six flip strata at 200
draws, and negative in the three tiered strata and in the 1,000-draw row ($-0.200$). «Strictly increasing the risk»: a claim that Meek already
derives leaves $\Gz$, $Z$ and your $\rval$ unchanged.
\end{armadilha}

\prompt[conceito=frases-proibidas,dificuldade=0.75]{Give the two Table 1 cells that contradict «a strong
and consistent predictor», and the stratum where both baselines score above the radius.}%
{$\tau_b(\rval) = -0.091$ in flip, cov 1.0, bw 0.00, and the confirmed null $0.032\ [-0.112, 0.168]$ in
flip, cov 0.5, bw 0.25 at 1,000 draws. In flip, cov 1.0, bw 0.25, SHD scores 0.582 and $|\Kset|$ 0.557
against 0.355 for the radius.}

\subsection*{8. The bibliography is broken in the compiled PDF}

\emph{Where.} \texttt{keel2024fuzzy} and \texttt{robust2026imperfect} have no author field, so they print
from their titles and are cited as «(kee, 2024)» and «(rob, 2025)»; \texttt{ejaz2025lges} prints the wrong
title and first name. \emph{Why.} These are visible on page~3 and page~12 of the PDF, and one of them puts a
researcher's work under a name that is not theirs. \emph{Fix.} From the arXiv records: arXiv:2405.08699 is
by Wenrui Li, Wei Zhang, Qinghao Zhang, Xuegong Zhang and Xiaowo Wang; arXiv:2511.06790 by Zidong Wang, Xi
Lin, Chuchao He and Xiaoguang Gao; arXiv:2506.22331 is «Less Greedy Equivalence Search», by Adiba Ejaz
and Elias Bareinboim. Your own bib comments already flag all three with «VERIFY».

\subsection*{9. Citations a referee will ask for}

\emph{Where.} Related work (\S3). \emph{Why.} The central word is «breakdown», the distance is evaluated on
adjustment, and the evaluation runs a tiered arm, so four literatures are one question away. \emph{Fix.}
Cite Masten and Poirier, «Inference on breakdown frontiers» (Quantitative Economics 11(1):41--111, 2020);
the adjustment identification distance (Henckel, Würtzen and Weichwald, UAI 2024) and the structural
intervention distance (Peters and Bühlmann, Neural Computation 27:771--799, 2015); Bang and Didelez on
tiered knowledge, both entries already in your bibliography and never cited; and b-LOAD
(\texttt{ahn2026bload}), also in the bibliography and never cited. Chapter~\ref{ch:where} adds the minimal
correction set.

\subsection*{10. The main text is eleven pages}

\emph{Where.} The conclusion ends on page~11 of the compiled draft. \emph{Why.} The ICLR 2027 author
guidelines say «At the time of submission, the main text should be 9 pages or fewer» and «Papers with main
text beyond the page limit will be desk-rejected». \emph{Fix.} Chapter~\ref{ch:merge}, where the cut is part
of the merge.

\textbf{Smaller items.} «What is required is model inclusion» in \S4.3 is necessary and not sufficient
(Exercise~\ref{ex:review-membership}). «All 133 CPDAGs at four nodes» in Appendix~B mixes two populations:
\texttt{rho\_breakdown/\allowbreak theory\_checks\_2026-09-16/\allowbreak check\_four\_node\_classes.py} counts 185 labelled
four-node classes, 126 of them with an undirected edge, and 133 is those 126 plus the 7 three-node classes
with one, which is what A.1 says. An unreachable failure has three conventions ($\infty$ in Definition~3, a
status in \S5.4, \mbox{\texttt{UNREACHED = -1}} in the code); pick one.

\naopasse[distancia-em-saltos]{Write out the upper bound $d(G, H) \le |\Kset_G \triangle \Kset_H|$ and name
the exact place where Lemma~2 is used in it. If you can only write the lower bound, reread item~3.}

\subsection*{Exercises}

\begin{exercicio}[metodo=counterexample,conceito=espaco]
\label{ex:review-membership}
Your \S4.3 says «What is required is model inclusion, $[G] \subseteq [\Chat]$, which the fixpoint test
enforces». Give the smallest counterexample showing that model inclusion alone does not make $G$ a knowledge
state, check it with the fixpoint test, and name the property it lacks.
\end{exercicio}
\begin{solucao}
Take $\Chat = A - B - C$ (three DAGs) and $G$ with $A\to B$ and $B - C$. Then $[G] = \{A\to B\to C\}$ is a
non-empty strict subset of $[\Chat]$, so model inclusion holds. But $\mathrm{Meek}(\Chat, \{A\to B\})$
applies R1 ($A$ and $C$ are non-adjacent) and returns $A\to B\to C \ne G$, so the fixpoint test fails. $G$
lacks maximal orientation: $B - C$ is undirected although every DAG of $[G]$ orients it $B\to C$. The test
enforces inclusion \emph{and} maximal orientation, and the second conjunct does the work.
\end{solucao}

\begin{exercicio}[metodo=diagnosis,conceito=frases-proibidas]
A referee writes: «The authors claim their radius ranks survival better than standard baselines, and that
asserting more knowledge makes results more fragile. Table 1 supports neither.» Using only your Table~1,
diagnose which part is right, which cells carry it, and the smallest true replacement for each sentence.
\end{exercicio}
\begin{solucao}
The first part is right; the second holds for the flip strata and fails for the tiered ones. The radius is not consistent ($-0.091$ in flip, cov 1.0, bw 0.00; the null
$0.032\ [-0.112, 0.168]$ at 1,000 draws) and does not dominate (0.582 and 0.557 against 0.355 in flip, cov
1.0, bw 0.25). A true replacement is the commented-out paragraph of \S7: the radius shares significance with
the baselines in most strata and does not dominate them. For fragility: $\tau_b(|\Kset|)$ is negative in the three tiered strata ($-0.423$, $-0.382$,
$-0.377$) but positive in five of the six flip strata at 200 draws, so the true sentence is «under
tiered, block-correlated errors, more asserted claims go with lower survival», without «strictly».
\end{solucao}

\section*{Chapter map}
\begin{center}
\begin{tikzpicture}[node distance=0.6cm and 1.2cm]
\node[cmapg] (right) {what is right: object,\\ fixpoint, discipline};
\node[cmapg, right=of right] (theory) {theory re-implemented:\\ no counterexample};
\node[cmapr, right=of theory] (text) {what the text says};
\node[cmap, right=of text] (pos) {Prop.~9 positioning;\\ unstated conjecture};
\node[cmapr, below=of right] (res) {prose against table;\\ bibliography; length};
\node[cmapc, below=of theory] (merge) {the merge,\\ Chapter~\ref{ch:merge}};
\node[cmap, below=of text] (cond) {closed form; conditions\\ dropped; two criteria};
\draw[cedge] (right) -- node[above]{kept} (theory);
\draw[cedge] (theory) -- node[above]{fixes lie in} (text);
\draw[cedge] (text) -- node[above]{positioning} (pos);
\draw[cedge] (text) -- node[right]{statements} (cond);
\draw[cedge] (res) -- node[above]{rewrite, cut} (merge);
\draw[cedge] (cond) -- node[above]{clauses} (merge);
\draw[cedge] (pos.south) to[out=-90, in=-25, looseness=0.55] node[below, pos=0.25]{one paragraph} (merge.south east);
\end{tikzpicture}
\end{center}
